\section{Locked half-suits, declaration timing, and termination}
\label{sec:locked}

Declaration is the only action in the dialect of \S\ref{sec:rules} that converts
information into score. This section states the one structural property of
declaration timing that follows from the rules alone, separates it from the
properties it does \emph{not} imply, gives the fitted rule \vfast{} uses, and
reports what is and is not known about termination.

\subsection{A half-suit held by one team cannot be taken back}

\begin{theorem}[Locked half-suits]
\label{thm:locked}
Suppose every card of half-suit $H$ is held by members of team $T$. Then no
member of the opposing team can ever legally ask for a card of $H$, so no card of
$H$ can change hands, and $H$ remains wholly held by $T$ for as long as it stays
undeclared. Moreover any declaration of $H$ by the opposing team is necessarily
incorrect and awards $H$ to $T$.
\end{theorem}

\begin{proof}
A legal ask for a card of $H$ requires the asker to hold another card of $H$, so
with all six cards held by members of $T$ every opposing player is void in $H$
and has no such ask, and a player holding no cards at all has no legal ask of any
kind. Cards move only through successful asks, and the only asks in $H$ available
to anyone are asks by members of $T$ directed at opponents, which are guaranteed
misses; ownership of $H$ is therefore frozen. For the second part, a declaration
must assign every card of $H$ to a teammate of the declarer, so an opposing
declarer would assign cards held by $T$ to the other team, which is false, and
$H$ is awarded to $T$.
\end{proof}

Theorem~\ref{thm:locked} does not assert that the probability of a correct claim
rises monotonically, and it does not: as a ratio of counts of consistent
assignments, both numerator and denominator shrink as constraints accumulate, and
a deduction elsewhere in the deal can eliminate more completions of the true
allocation than of a rival. What is monotone is the support: the allocations of
$H$ extendable to a consistent deal are non-increasing.

\begin{corollary}[Waiting carries no ownership risk]
\label{cor:wait}
For a locked half-suit the ``steal'' term in the declaration trade-off is
identically zero: $H$ can leave the undeclared state only through a declaration
by $T$, or one by the opposing team, which by Theorem~\ref{thm:locked} awards $H$
to $T$ in any case. Declaring immediately therefore cannot protect ownership; the
remaining reasons to declare are the point itself and the positional consequences
of removing six cards from play.
\end{corollary}

Corollary~\ref{cor:wait} concerns ownership only, which matters for
\S\ref{sec:termination}.

\begin{observation}[Ownership freezes; allocation information need not]
\label{obs:frozeninfo}
Let $H$ be a half-suit all of whose cards sit with team $T$. Two statements hold
simultaneously and should not be merged. First, ownership of $H$ is frozen by
Theorem~\ref{thm:locked}. Second, asks inside $H$ remain \emph{legal} for members
of $T$; such an ask transfers nothing, every target being void, but it is a
public event and still carries the certificate \eqref{eq:disj}, which constrains
which teammate holds which card. Allocation information about $H$ can therefore
continue to change while ownership cannot, deliberately at the price of a turn
and incidentally as capacities tighten elsewhere through \eqref{eq:capacity}.
Ownership monotonicity alone therefore does not establish that a position with
only locked half-suits is informationally frozen, and does not prove universal
non-termination.
\end{observation}

Develin \cite{develin} records the same channel from the human side, and
\vfast{} does not exploit it deliberately (\S\ref{sec:limitations}). Two forces
therefore act on declaration timing in opposite directions --- a retained locked
half-suit keeps six cards absorbing capacity mass in $U$, while a hand of locked
cards licenses only asks certain to fail --- and which dominates is a
quantitative question about a position and an opponent population rather than one
settled by Theorem~\ref{thm:locked}. \S\ref{app:human} develops both sides.

\subsection{A value-based declaration rule}
\label{sec:stopping}

At each declaration opportunity \vfast{} chooses between converting a live
half-suit into a scored point and leaving it live, comparing two estimates of
position value rather than testing a confidence threshold. Let $V(\cdot)$ be the
fitted linear evaluation of \S\ref{sec:value}, in units of final half-suit
differential scaled to $[-1,1]$; $s$ the current position, $\Delta$ the declaring
team's half-suit differential, $H$ the candidate half-suit, $A$ the named
allocation the policy would declare, and $\hat q = \hat\alpha(A)$ the Fast
estimate \eqref{eq:q2} that $A$ is correct. The rule declares when
\begin{equation}
  \hat q\, V\bigl(s \ominus H,\, \Delta{+}1\bigr)
  + (1 - \hat q)\, V\bigl(s \ominus H,\, \Delta{-}1\bigr)
  \;>\; V(s) + \varepsilon,
  \label{eq:stop}
\end{equation}
where $s \ominus H$ denotes the position with $H$ removed and $\varepsilon$ is a
fitted margin, one of the \numParams{} fitted coordinates of \S\ref{sec:fitting}.

Both sides of \eqref{eq:stop} are values of a fitted linear model evaluated on
features produced by an approximate transition, so it compares a fitted value
estimate under an approximate feature transition; it is not a solution to a
Bellman equation, and we make no optimality claim for it. The approximation is in
$s \ominus H$, where the implementation updates aggregate team-level quantities
only: it removes the half-suit's contributions to expected control, sharpened
control, the locked differential and the contested mass, removes its unresolved
cards from $U$, and reduces the declaring team's card count by six. That
card-count reduction is applied on \emph{both} the correct and the incorrect
branch, although on an incorrect declaration some of those six cards were held by
opponents and the true post-declaration counts differ, and player-specific
features are not re-derived on either branch. The resulting feature vector is an
approximate summary of the successor position rather than the successor position
itself.

What distinguishes \eqref{eq:stop} from a threshold on $\hat q$ is which terms
survive. The immediate expected score change from declaring is $2\hat q - 1$; but
a value function fitted on positions in which the half-suit is still live already
credits the team's expected control of it, so that term is largely cancelled by
the control and locked-differential contributions removed in $s \ominus H$, and
what remains to decide the comparison are the positional terms
Corollary~\ref{cor:wait} identifies --- the information the six unresolved slots
absorb against the turn efficiency of a lighter hand. In the limiting case
$\hat q = 1$ with $H$ locked, the score term vanishes entirely and
\eqref{eq:stop} is a comparison between exactly those two effects. This is how
the rule is constructed; the ablation below is what the construction is worth,
and it resolves nothing.

The frozen-policy ablation does not isolate a benefit from this formulation.
Replacing \eqref{eq:stop} with a fixed confidence threshold changed the win rate
by \numAblStop{} points, 95\% paired CI \numAblStopCI{} (\S\ref{sec:ablations}),
and moving the declaration confidence threshold of the shipped configuration
between 0.80 and 0.99 also changes nothing the experiment can resolve. One reason
the achievable effect is small is the resolution of the model class $V$ is drawn
from: a linear function of \numValueFeats{} summary features caps what any rule
built on differences of $V$ can extract. The only recorded fit in that class is
E14, at $R^2 = \numValueRSq$ in sample on \numValueRows{} self-play decision
points; its coefficients are not the ones compiled into the shipping policy
(\S\ref{sec:value}, \S\ref{app:params}), so that figure describes the class and
not the deployed vector.

Measured over the primary head-to-head bank of \S\ref{sec:results-h2h}, \vfast{}
holds a half-suit that has become locked for \numLockHoldFour{} public events
before declaring it, against \numLockHoldThree{} for \vthree{}. Three
qualifications apply and none can be removed: the moment a half-suit becomes
locked is computed from ground truth after the fact and is never available to
either policy; both figures are conditioned on the declaration being correct; and
the two policies declare correctly at very different rates
(\numVsVthreeDeclAcc{}\% and \numVthreeDeclAccSame{}\% on these same games), so
the averages are taken over differently selected subsets. This is measured
behaviour, not a theorem: the qualifications leave the magnitude of the
difference unusable, and only the sign survives all three
(\S\ref{app:human}). The direction is at least consistent with the two
estimators: \vthree{} scores a candidate allocation by a geometric mean of
per-card marginals, which needs many more public events to clear a threshold than
the joint estimate \eqref{eq:q2} does, so the earlier declaration is what a
sharper allocation probability would be expected to produce. That is an account,
not a measurement; no experiment here isolates it.

\subsection{Termination}
\label{sec:termination}

Nothing in the rules of the dialect studied compels a declaration. Consider a
position in which every live half-suit is wholly owned by one team or the other,
but at least one team cannot resolve which of its own members holds which card:
by Theorem~\ref{thm:locked} card movement ceases, every player being void in
every live half-suit owned by the opposing team, so every legal ask is a
guaranteed miss. Whether the position is absorbing depends on whether further
allocation information can arrive, which by Observation~\ref{obs:frozeninfo} is
not settled by Theorem~\ref{thm:locked} alone, those misses being legal, public
and certificate-carrying. A policy whose declaration rule is sensitive to the
certificates may eventually declare from such a position; one whose rule is not
may not.

\begin{observation}[Observed non-termination in mirror self-play]
\label{obs:cycle}
An earlier fitted configuration of the policy, played against a copy of itself
with no forcing rule, failed to terminate within the action cap in
\numMirrorCycle\% of games at the 400-ask cap of \S\ref{sec:rules}. Raising the
cap to 1000 asks left the figure unchanged, which distinguishes a cycle from slow
convergence (E11, 150 deals in two seat orientations, seed 31). This is a
measured property of that fitted configuration on that corpus, not a theorem
about the rules dialect.
\end{observation}

The remedy adopted is in the policy. A test that forces a declaration as soon as
no productive ask remains also fires in ordinary early positions, where a player
holding only cards their own team already owns has no productive ask but ample
reason to wait; wiring it in reduced the mirror cap-hit rate to
\numDeadAskCycle\% but cost \numDeadAskCost{} win-rate points against \vthree{}
(E11). The shipped alternative is a graduated escalation on the public event
count: below a forcing horizon \eqref{eq:stop} decides, past it the policy
declares any half-suit whose estimated allocation probability exceeds one half,
and past a second horizon its best candidate whatever the estimate, an
undeclared half-suit scoring nothing. Under this rule every game terminated
through play and the action-limit rate is \numLimitRate{} in E1, E3, E7, E10 and
E13; the exception is the deliberately unforced mirror configuration of
Observation~\ref{obs:cycle}, measured precisely because it does reach the cap.
Both horizons are hand-set constants chosen during development, not the
output of any analysis in this paper, and no sensitivity study over them is
reported.

A general termination result would require more than we have: either a
constructive reachable position, with a strategy profile, in which the full
information state as well as the physical position recurs, so that play is
provably periodic --- Observation~\ref{obs:frozeninfo} is the obstacle, since the
certificate stream from legal misses must also be shown to add nothing --- or a
proof over the joint dynamics of the full information state and the
action-selection rule of the policies involved. We have neither, and
Observation~\ref{obs:cycle} is evidence only that non-termination is reachable in
practice for some fitted configurations. The rules literature
already records a forced-claims proposal for tournament play \cite{pagat}, with
which these measurements are consistent.
